
% Para una visualizacion correcta, generar el PDF
% Ni el DVI ni el PS se visualizan bien

% Elegir el estilo que se desee, hay cientos en la red

\documentclass{beamer}
\usepackage{beamerthemeshadow}
\usepackage[spanish]{babel}
\usepackage[utf8]{inputenc}
%\usepackage[latin1]{inputenc}

\begin{document}
\title{Laboratorio virtual para la simulación y análisis de paradigmas de toma de decisiones humanas mediante agentes inteligentes}
\subtitle{Fase 1: generación automática de modelos de decisión \\
Grado en Ingeniería Informática \\
Universidade de Santiago de Compostela}
\author{Autor: Pablo Pazos Parada}
\institute{Tutor: Eduardo Manuel Sanchez Vila}
\date{\today}

\begin{frame}
\titlepage
\end{frame}

\begin{frame}
\frametitle{Tabla de contenidos}\tableofcontents
\end{frame}

\section{Problema y objetivo}
\begin{frame}
\frametitle{Problema}
Convertir literatura y descripciones de paradigmas de decisión en modelos
ejecutables sigue siendo una tarea manual, lenta y difícil de auditar.
\end{frame}

\begin{frame}
\frametitle{Objetivos}
\begin{itemize}
\item Investigar paradigmas científicos relevantes.
\item Formalizarlos matemáticamente.
\item Adaptarlos a un entorno de simulación.
\item Generar código Python probado.
\end{itemize}
\end{frame}

\section{Arquitectura}
\begin{frame}
\frametitle{Pipeline}
\begin{center}
Researcher $\rightarrow$ Formalizer $\rightarrow$ Reasoner $\rightarrow$ Builder
\end{center}
\end{frame}

\begin{frame}
\frametitle{Memoria}
\begin{itemize}
\item MinIO: artefactos.
\item PostgreSQL: metadatos y ciclo de vida.
\item Neo4j: entidades y relaciones.
\item Qdrant: recuperación densa y BM25.
\end{itemize}
\end{frame}

\section{Resultados}
\begin{frame}
\frametitle{Salida}
Modelos Python autocontenidos compatibles con:
\begin{itemize}
\item \texttt{decide(perception)}
\item \texttt{update(action, reward, new\_perception)}
\item \texttt{get\_state()}
\end{itemize}
\end{frame}

\begin{frame}
\frametitle{Conclusión}
\begin{itemize}
\item Pipeline trazable y revisable.
\item Memoria persistente para reutilizar conocimiento aceptado.
\item Validación mediante estructura, revisión y tests.
\end{itemize}
\end{frame}

\end{document}
